\documentclass[11pt,a4paper]{article}
\usepackage[italian]{babel}
\usepackage[margin=2.5cm]{geometry}
\usepackage{parskip}
\usepackage{amsmath, amssymb, amsthm}
\usepackage{booktabs}
\usepackage{array}
\usepackage{xcolor}
\usepackage{needspace}
\usepackage{enumitem}
\usepackage{listings}
\usepackage{tikz}
\usepackage[most]{tcolorbox}
\usepackage[hidelinks]{hyperref}
\usetikzlibrary{decorations.pathreplacing, shapes.geometric, arrows.meta, positioning}

% Niente vedove/orfane: i paragrafi non lasciano righe isolate a fine/inizio pagina
\widowpenalty=10000
\clubpenalty=10000

% Elenchi più compatti
\setlist{itemsep=2pt, parsep=0pt, topsep=4pt, partopsep=0pt}

% --- Riquadri con bordo nero, senza numero (si spezzano tra due pagine) ---
% Uso: \begin{definizione} ... \end{definizione}
%      \begin{teorema}[Nome] ... \end{teorema}  (il nome è facoltativo)
\tcbset{nota/.style={enhanced jigsaw, breakable, colback=white,
  colframe=black, boxrule=0.5pt, sharp corners}}
\NewTColorBox{definizione}{}{nota,
  before upper={\textbf{Definizione.}\ }}
\NewTColorBox{teorema}{o}{nota,
  before upper={\textbf{Teorema\IfValueT{#1}{ (#1)}.}\ }}
\NewTColorBox{proposizione}{o}{nota,
  before upper={\textbf{Proposizione\IfValueT{#1}{ (#1)}.}\ }}
\NewTColorBox{proprieta}{o}{nota,
  before upper={\textbf{Proprietà\IfValueT{#1}{ (#1)}.}\ }}
\NewTColorBox{assioma}{o}{nota, boxrule=1pt,
  before upper={\textbf{Assioma\IfValueT{#1}{ (#1)}.}\ }}

% --- Ambienti semplici (senza riquadro) ---
% Il titolo sta in cima al blocco, anche se il contenuto inizia con un riquadro o
% con codice; e non resta mai da solo in fondo a una pagina.
% Uso: \begin{esempio}[Titolo facoltativo] ... \end{esempio}
\newcommand{\newrunin}[2]{%
  \NewDocumentEnvironment{#1}{o}{%
    \par\Needspace{4\baselineskip}\addvspace{\medskipamount}\noindent
    \textbf{#2}\IfValueT{##1}{ (##1)}\textbf{.}\ \ignorespaces}%
  {\par\addvspace{\medskipamount}}}
\newrunin{esempio}{Esempio}
\newrunin{esercizio}{Esercizio}
\NewTColorBox{osservazione}{}{nota,
  before upper={\textbf{Osservazione.}\ }}

% V e F nelle tavole di verità
\newcommand{\V}{\textsf{V}}
\newcommand{\F}{\textsf{F}}

% --- Codice C ---
% Uso: \begin{codice} ... \end{codice}      codice C numerato
%      \begin{comando} ... \end{comando}    comandi da terminale
%      \begin{risultato} ... \end{risultato} output di un programma
%      \lstinline|printf("ciao");|           codice dentro al testo
\lstdefinestyle{c}{
  language=C,
  basicstyle=\ttfamily\small,
  keywordstyle=\color{blue}\bfseries,
  commentstyle=\color{gray}\itshape,
  stringstyle=\color{red!70!black},
  numbers=left,
  numberstyle=\tiny\color{gray},
  numbersep=8pt,
  columns=flexible,
  keepspaces=true,
  breaklines=true,
  showstringspaces=false,
  tabsize=4,
  morekeywords={include, define},
  % lettere accentate nei commenti
  literate={à}{{\`a}}1 {è}{{\`e}}1 {é}{{\'e}}1 {ì}{{\`i}}1
           {ò}{{\`o}}1 {ù}{{\`u}}1 {È}{{\`E}}1 {À}{{\`A}}1
}
\lstset{style=c}
\newtcblisting{codice}{listing only, nota, left=6mm,
  listing options={style=c}}
\newtcblisting{comando}{listing only, nota,
  listing options={basicstyle=\ttfamily\small, numbers=none, language={},
    columns=flexible, keepspaces=true, breaklines=true}}
\newtcblisting{risultato}{listing only, enhanced jigsaw, breakable,
  colback=black!5, colframe=black, boxrule=0.5pt, sharp corners,
  listing options={basicstyle=\ttfamily\small, numbers=none, language={},
    columns=flexible, keepspaces=true, breaklines=true}}

% --- Diagrammi di flusso (simboli standard) ---
\tikzset{
  terminale/.style = {ellipse, draw, minimum width=2.5cm, minimum height=0.9cm},
  io/.style        = {trapezium, trapezium left angle=70, trapezium right angle=110,
                      draw, minimum width=2.5cm, minimum height=0.9cm},
  processo/.style  = {rectangle, draw, minimum width=2.5cm, minimum height=0.9cm},
  decisione/.style = {diamond, draw, aspect=2, minimum width=2.5cm},
  freccia/.style   = {thick, -{Stealth}}
}

\title{Tipi di dato}
\author{Marco Cerbella}
\date{Lezione 8 -- 7 ottobre 2026}

\begin{document}
\maketitle

\section{Tipi e oggetti}

I programmi devono memorizzare ed elaborare dati di tipo diverso, come interi e
numeri in virgola mobile, in modi diversi: il compilatore deve quindi sapere che
tipo di dato rappresenta un valore.

\begin{definizione}
In C un \emph{oggetto} è una zona di memoria il cui contenuto può rappresentare
valori. Gli oggetti che hanno un nome sono detti \emph{variabili}. Il
\emph{tipo} di un oggetto determina:
\begin{itemize}
    \item quanto spazio occupa in memoria;
    \item i valori che la variabile può assumere;
    \item le operazioni che si possono eseguire su di essa.
\end{itemize}
\end{definizione}

\subsection{Tipi in C}

\begin{itemize}
    \item \textbf{Tipi di base}: interi (standard ed estesi) e floating-point
          (reali e complessi).
    \item \textbf{Tipi enumerati}.
    \item Il tipo \lstinline|void|.
    \item \textbf{Tipi derivati}: puntatori, array, strutture, unioni, funzioni.
\end{itemize}

I tipi di base e i tipi enumerati formano i \emph{tipi aritmetici}; i tipi
aritmetici e i puntatori formano i \emph{tipi scalari}; array e strutture sono
detti \emph{tipi aggregati}. Un \emph{tipo funzione} descrive l'interfaccia di
una funzione: il tipo del valore restituito ed eventualmente i tipi dei
parametri.

\section{Interi}

I tipi interi con segno sono cinque (molti hanno sinonimi):
\begin{center}
\begin{tabular}{ll}
\toprule
\textbf{tipo} & \textbf{sinonimi} \\
\midrule
\lstinline|signed char| & \\
\lstinline|int| & \lstinline|signed|, \lstinline|signed int| \\
\lstinline|short| & \lstinline|short int|, \lstinline|signed short|, \lstinline|signed short int| \\
\lstinline|long| & \lstinline|long int|, \lstinline|signed long|, \lstinline|signed long int| \\
\lstinline|long long| (C99) & \lstinline|long long int|, \lstinline|signed long long|, \lstinline|signed long long int| \\
\bottomrule
\end{tabular}
\end{center}

C definisce solo le dimensioni minime: \lstinline|short| occupa almeno 2 byte,
\lstinline|long| almeno 4 e \lstinline|long long| almeno 8. I tipi possono essere
più grandi, ma devono rispettare l'ordine
\[
  \texttt{sizeof(short)} \leq \texttt{sizeof(int)} \leq \texttt{sizeof(long)}
  \leq \texttt{sizeof(long long)} .
\]

\subsection{Booleani}

In C non esistono \emph{true} e \emph{false}: $0$ è falso e qualsiasi valore
diverso da $0$ (ad esempio $1$, $-25$, $123456$) è vero.

\begin{codice}
if (3)                      if (0)
    printf("YES\n");            printf("YES\n");
else                        else
    printf("NO\n");             printf("NO\n");
\end{codice}

Il primo \lstinline|if| stampa \texttt{YES}, il secondo \texttt{NO}.

C99 ha introdotto il tipo intero senza segno \lstinline|_Bool|: \emph{vero} è
codificato come $1$, \emph{falso} come $0$. Includendo l'header
\lstinline|stdbool.h| si possono usare anche \lstinline|bool|, \lstinline|true| e
\lstinline|false|: \lstinline|bool| è sinonimo di \lstinline|_Bool| e
\lstinline|true|, \lstinline|false| sono costanti uguali a $1$ e $0$.

\subsection{Caratteri}

Anche \lstinline|char| è un tipo intero standard, ma il nome \lstinline|char| è
sinonimo di \lstinline|signed char| oppure di \lstinline|unsigned char| a seconda
del compilatore. Occupa 1 byte e i caratteri corrispondono ai codici della
tabella ASCII.

Con le variabili \lstinline|char| si può fare aritmetica: sta al programmatore
decidere se il numero contenuto è un codice di carattere o altro.

\begin{codice}
char ch = 'A';       // variabile di tipo char
printf("The character %c has the character code %d.\n", ch, ch);
printf("%c", ch + 1);
\end{codice}
\begin{risultato}
The character A has the character code 65.
B
\end{risultato}

\section{Sistemi di numerazione}

\begin{definizione}
Il sistema \emph{binario} (base 2) rappresenta i numeri con due soli simboli,
$0$ e $1$; ogni cifra è un \emph{bit}. Quasi tutti i calcolatori lo usano
internamente perché è semplice da realizzare con circuiti a porte logiche.
\end{definizione}

\begin{definizione}
Il sistema \emph{esadecimale} (base 16) usa sedici simboli: le cifre $0$--$9$ per
i valori da zero a nove e le lettere \texttt{A, B, C, D, E, F} (o minuscole) per
i valori da dieci a quindici. Ogni cifra esadecimale corrisponde a 4 bit (un
\emph{nibble}, metà di un byte da 8 bit): è una rappresentazione più
leggibile dei valori binari. Il sistema \emph{ottale} (base 8) usa le cifre
$0$--$7$.
\end{definizione}

\begin{center}
\begin{tabular}{rcc @{\qquad} rcc}
\toprule
\textbf{Bin.} & \textbf{Esa.} & \textbf{Dec.} &
\textbf{Bin.} & \textbf{Esa.} & \textbf{Dec.} \\
\midrule
0    & 0 & 0  & 1000 & 8 & 8 \\
1    & 1 & 1  & 1001 & 9 & 9 \\
10   & 2 & 2  & 1010 & A & 10 \\
11   & 3 & 3  & 1011 & B & 11 \\
100  & 4 & 4  & 1100 & C & 12 \\
101  & 5 & 5  & 1101 & D & 13 \\
110  & 6 & 6  & 1110 & E & 14 \\
111  & 7 & 7  & 1111 & F & 15 \\
\bottomrule
\end{tabular}
\end{center}

\subsection{Da binario/esadecimale a decimale}

Un numero con cifre $d_n \dots d_1 d_0$ in base $b$ vale
$d_n b^n + d_{n-1} b^{n-1} + \dots + d_1 b^1 + d_0 b^0$ (in base 2:
$d_n 2^n + \dots + d_1 2^1 + d_0 2^0 = N_{10}$).

\begin{esempio}
$1001_2 = 1 \cdot 2^3 + 0 \cdot 2^2 + 0 \cdot 2^1 + 1 \cdot 2^0 = 9$ \quad e \quad
$\texttt{002B}_{16} = 2 \cdot 16 + 11 = 43$.
\end{esempio}

\subsection{Da decimale a binario/esadecimale}

Si divide ripetutamente il numero per la base e si leggono i resti
\emph{dal basso verso l'alto}. Per l'ottale si usa lo stesso algoritmo con base 8.

\begin{esempio}
$156$ in binario (divisioni per 2): i quozienti successivi sono
$78, 39, 19, 9, 4, 2, 1, 0$ e i resti $0, 0, 1, 1, 1, 0, 0, 1$. Letti dal basso:
$156 = 10011100_2$.

$1565$ in esadecimale (divisioni per 16): $1565 = 97 \cdot 16 + 13$ (resto
$13 = \texttt{d}$), $97 = 6 \cdot 16 + 1$ (resto $1$), $6 = 0 \cdot 16 + 6$
(resto $6$). Quindi $1565 = \texttt{61d}_{16}$.
\end{esempio}

Con $N$ bit si rappresentano i numeri da $0$ a $2^N - 1$; con 8 bit (un byte)
da $0$ a $255$ (da \texttt{0000 0000} a \texttt{1111 1111}). Servono però anche i
numeri negativi, e per questo esistono rappresentazioni diverse: segno e modulo,
complemento a due.

\section{Rappresentazione degli interi in memoria}

\subsection{Segno e modulo}

Un bit (di solito il più a sinistra) indica il segno: $0$ per i positivi, $1$ per
i negativi; i bit restanti rappresentano il modulo. Ad esempio \texttt{1001 1000}
vale $-24$ (mentre se si usasse solo per numeri positivi varrebbe $152$).

Con $n$ bit si rappresentano i valori da $-(2^{n-1} - 1)$ a $2^{n-1} - 1$, e
$\pm 0$: con 8 bit da $-127$ a $+127$.

\begin{osservazione}
Problema: lo zero ha due rappresentazioni, \texttt{0000 0000} ($+0$) e
\texttt{1000 0000} ($-0$). Una soluzione è il complemento a due.
\end{osservazione}

\subsection{Complemento a due}

\begin{center}
\begin{tabular}{ccc}
\toprule
\textbf{Valore binario} & \textbf{Complemento a due} & \textbf{Senza segno} \\
\midrule
\texttt{00000000} & $0$    & $0$   \\
\texttt{00000001} & $1$    & $1$   \\
$\vdots$ & $\vdots$ & $\vdots$ \\
\texttt{01111110} & $126$  & $126$ \\
\texttt{01111111} & $127$  & $127$ \\
\texttt{10000000} & $-128$ & $128$ \\
\texttt{10000001} & $-127$ & $129$ \\
\texttt{10000010} & $-126$ & $130$ \\
$\vdots$ & $\vdots$ & $\vdots$ \\
\texttt{11111110} & $-2$   & $254$ \\
\texttt{11111111} & $-1$   & $255$ \\
\bottomrule
\end{tabular}
\end{center}

Nel complemento a due esiste un solo zero, \texttt{00000000}. Per cambiare segno a
un numero (positivo o negativo) si \textbf{invertono tutti i bit e poi si somma 1}.

\begin{esempio}
\begin{itemize}
    \item \texttt{0000 0101} (valore $5$); invertendo: \texttt{1111 1010};
          sommando $1$: \texttt{1111 1011}, che vale $-5$.
    \item Da \texttt{1111 1011}: invertendo \texttt{0000 0100}, sommando $1$
          \texttt{0000 0101}, cioè di nuovo $5$.
\end{itemize}
\end{esempio}

Con $n$ bit si rappresentano i valori da $-2^{n-1}$ a $2^{n-1} - 1$: non c'è
``$-0$'', quindi si rappresenta un numero negativo in più. Con 8 bit, da $-128$
(\texttt{1000 0000}) a $+127$ (\texttt{0111 1111}).

\begin{osservazione}
La regola ``inverti e somma 1'' non funziona per $-128$, perché $128$ non è
rappresentabile con 8 bit in complemento a due. D'ora in poi si usa il complemento
a due.
\end{osservazione}

\subsection{Endianness}

\begin{definizione}
L'\emph{endianness} indica l'ordine usato per interpretare numericamente una
sequenza di byte della memoria come una parola più grande (e l'ordine di
trasmissione dei byte su un collegamento digitale). Le parole possono essere in
formato \emph{big-endian} o \emph{little-endian}, a seconda che si parta
dall'estremità più significativa (big end) o meno significativa (little end).
\end{definizione}

Esempi: i mainframe IBM z/Architecture e i Motorola 68000 usano big-endian, i
processori Intel x86 little-endian.

La memoria è un array di byte; ogni byte ha il suo \emph{indirizzo logico}
(un numero positivo), cioè l'indirizzo al quale un elemento appare risiedere dal
punto di vista di un programma in esecuzione. Per architetture a 64 bit il limite
superiore è $2^{64} - 1$.

\begin{esempio}
Il numero $2$ scritto su 4 byte, a partire dall'indirizzo $1500$:
\begin{center}
\begin{tabular}{lcccc}
\toprule
 & 1500 & 1501 & 1502 & 1503 \\
\midrule
big-endian    & \texttt{00} & \texttt{00} & \texttt{00} & \texttt{02} \\
little-endian & \texttt{02} & \texttt{00} & \texttt{00} & \texttt{00} \\
\bottomrule
\end{tabular}
\end{center}
\end{esempio}

\begin{esempio}
Sulla stringa di 4 bit \texttt{0010}, in complemento a due: in big-endian vale $2$;
in little-endian si leggono i bit in ordine inverso (\texttt{0100}) e vale $4$.
Sulla stringa \texttt{1010}: in big-endian vale $-6$, in little-endian
(\texttt{0101}) vale $5$.
\end{esempio}

\section{Dimensione e limiti dei tipi}

L'operatore \lstinline|sizeof| restituisce la dimensione esatta in byte di un tipo
o di una variabile: \lstinline|sizeof(tipo)| oppure \lstinline|sizeof espressione|
(in questo caso conta il tipo dell'espressione).

\begin{codice}
int iIndex;
iIndex = 1000;
// sizeof(int) e sizeof(iIndex) valgono entrambi 4
\end{codice}

I limiti dei tipi interi per il proprio compilatore sono nell'header
\lstinline|limits.h|, che definisce macro come \lstinline|INT_MIN|,
\lstinline|INT_MAX|, \lstinline|UINT_MAX|, \lstinline|CHAR_MIN|,
\lstinline|CHAR_MAX|.

\begin{codice}
#include <stdio.h>
#include <limits.h>

int main() {
    printf(" char %zu %d %d\n", sizeof(char), CHAR_MIN, CHAR_MAX);
    printf(" int %zu %d %d\n", sizeof(int), INT_MIN, INT_MAX);
    return 0;
}
\end{codice}
\begin{risultato}
 char 1 -128 127
 int 4 -2147483648 2147483647
\end{risultato}

\begin{osservazione}
Il risultato di \lstinline|sizeof| ha tipo \lstinline|size_t|: per stamparlo è
corretto \lstinline|%zu| (nelle slide compare \lstinline|%d|). I valori stampati
dipendono dalla piattaforma (qui un caso tipico).
\end{osservazione}

\section{Numeri in virgola mobile}

Per rappresentare i non interi, con la virgola in qualsiasi posizione, C ha tre
tipi floating-point standard:
\begin{itemize}
    \item \lstinline|float|: precisione singola;
    \item \lstinline|double|: precisione doppia;
    \item \lstinline|long double|: precisione estesa.
\end{itemize}

\subsection{Precisione}

Un valore in virgola mobile si memorizza solo con una \emph{precisione} limitata,
determinata dal formato binario e dalla memoria usata. La precisione si esprime in
cifre significative: ad esempio con sei cifre significative la conversione in un
numero decimale di sei cifre restituisce le sei cifre originali. La posizione della
virgola e gli zeri iniziali o finali non contano: $123\,456\,000$ e
$0{,}00123456$ si possono entrambi memorizzare in un tipo con precisione a sei
cifre.

In C le operazioni aritmetiche tra numeri in virgola mobile sono eseguite
internamente con precisione double o superiore.

\begin{codice}
float height = 1.2345, width = 2.3456;   // precisione singola
double area = height * width;            // il calcolo è fatto in double
\end{codice}

Se il risultato è assegnato a una variabile \lstinline|float|, viene arrotondato
se necessario. L'header \lstinline|float.h| definisce macro come
\lstinline|FLT_MIN|, \lstinline|FLT_MAX| e \lstinline|FLT_DIG|, che indicano
intervallo di valori e precisione del tipo \lstinline|float|.

\subsection{Notazione E e formato IEEE 754}

La \emph{notazione E} è la rappresentazione testuale della notazione scientifica:
\texttt{1.234e+56} significa $1{,}234 \times 10^{56}$.

Nel formato \emph{IEEE 754} ogni numero finito è descritto da tre interi: un segno
$s$ (zero o uno), un \emph{significando} (o mantissa) $c$ e un esponente $q$. Il
valore è
\[
  (-1)^s \times c \times b^q ,
\]
dove $b$ è la base (2 o 10). Ad esempio con base 10, segno $1$ (negativo),
significando $12345$ ed esponente $-3$ il valore è $-12{,}345$.

Il formato 754 a \emph{precisione singola} usa: 1 bit di segno, 8 bit di esponente,
23 bit di significando. Il valore è
\[
  (-1)^{\text{segno}} \times \Bigl( 1 + \sum_{i=1}^{23} b_{23-i}\, 2^{-i} \Bigr)
  \times 2^{(e - 127)} .
\]

\begin{esempio}
Con segno $0$, esponente \texttt{01111100} ($e = 124$) e frazione
\texttt{010\dots0} (cioè $b_{21} = 1$, tutti gli altri bit nulli) si ha
$1 + 2^{-2} = 1{,}25$ e $2^{124 - 127} = 2^{-3}$, quindi
\[
  \text{valore} = (+1) \times 1{,}25 \times 2^{-3} = +0{,}15625 .
\]
\end{esempio}

Con la \emph{precisione doppia} (\lstinline|double|) si usano 1 bit di segno,
11 di esponente e 52 di frazione; la \emph{precisione estesa}
(\lstinline|long double|, 80 bit) ha 1 bit di segno, 15 di esponente, 1 bit di
parte intera e 63 di frazione.

\subsection{Errori di arrotondamento}

L'errore di arrotondamento è intrinseco nel calcolo in virgola mobile.

\begin{codice}
#include <stdio.h>

int main() {
    int a = 16777217;
    float b = a;
    printf("%f\n", b);
}
\end{codice}
\begin{risultato}
16777216.000000
\end{risultato}

Il valore $16777217 = 2^{24} + 1$ non è rappresentabile esattamente con un
\lstinline|float|. Allo stesso modo conviene tenere in \lstinline|double| anche
il risultato di un prodotto come \lstinline|height * width|: salvarlo in un
\lstinline|float| fa perdere precisione.

\begin{osservazione}
Il modo più semplice per evitare di accumulare errori è usare numeri in virgola
mobile ad alta precisione, cioè \lstinline|double| al posto di \lstinline|float|.
Sulle CPU moderne non c'è (quasi) penalizzazione di tempo, ma un \lstinline|double|
occupa il doppio della memoria.
\end{osservazione}

\section{Tipi enumerati}

Le \emph{enumerazioni} sono tipi interi definiti dal programmatore. La definizione
inizia con la parola chiave \lstinline|enum|, eventualmente seguita da un
identificatore, e contiene l'elenco dei possibili valori con un nome per ciascuno:
\begin{codice}
enum [identificatore] { lista-enumeratori };
\end{codice}

\begin{esempio}
\begin{codice}
enum color { black, red, green, yellow, blue, white = 7, gray };
\end{codice}
Le costanti elencate valgono $0, 1, 2, 3, 4, 7, 8$ (se non si assegna un valore,
si continua dal precedente più uno).
\end{esempio}

\begin{codice}
enum color fgColor = blue,          // due variabili
           bgColor = yellow;        // di tipo enum color

void setFgColor(enum color fgc);    // funzione con parametro enum color
\end{codice}

Con le variabili di tipo enumerato si può fare aritmetica ordinaria
(\lstinline|red + red| vale $2$). Costanti diverse della stessa enumerazione
possono avere lo stesso valore:
\begin{codice}
enum signals { OFF, ON, STOP = 0, GO = 1, CLOSED = 0, OPEN = 1 };
\end{codice}

\begin{esempio}
\begin{codice}
#include <stdio.h>

enum week { sunday, monday, tuesday, wednesday, thursday, friday, saturday };

int main() {
    enum week today;
    today = wednesday;
    printf("Day %d", today + 1);
    return 0;
}
\end{codice}
\begin{risultato}
Day 4
\end{risultato}
\end{esempio}

\begin{osservazione}
Conviene usare sempre un'enumerazione quando una variabile (in particolare un
parametro di funzione) può assumere solo uno tra pochi valori possibili: si evitano
errori dovuti a costanti non valide, si documentano i valori ammessi ed è più
mnemonico degli interi.
\end{osservazione}

\section{Il tipo void}

Il tipo \lstinline|void| indica che non è disponibile alcun valore. Non si possono
dichiarare variabili o costanti di questo tipo, ma si usa:
\begin{itemize}
    \item nelle dichiarazioni di funzioni;
    \item nelle espressioni \lstinline|void|;
    \item nei puntatori a \lstinline|void| (si vedranno con i puntatori).
\end{itemize}

Una funzione che non restituisce nulla ha tipo \lstinline|void|:
\lstinline|void error(int a) {}|. \lstinline|void| nella lista dei parametri di un
prototipo indica che la funzione non ha parametri:
\lstinline|void printMenu(void) {}|. Il compilatore segnala un errore se si
chiama \lstinline|printMenu(3)|.

\end{document}
